% Sampling-distribution-CLT.tex
% Topic: Sampling Distribution of the Mean and the Central Limit Theorem

Suppose a population has mean $\mu$ and variance $\sigma^2$. We draw $n$
independent observations $X_1, X_2, \ldots, X_n$ from this population and
form the \textbf{sample mean}
\[
\bar{X} = \frac{1}{n}\sum_{i=1}^{n} X_i.
\]
Each $X_i$ is a random variable with the same distribution as the
population, and $\bar{X}$ is a random variable built from all of them.
Before any data are collected, $\bar{X}$ could take many values; its
distribution is called the \textbf{sampling distribution of the mean}.

\begin{remark}[Notation: $\bar{X}$ versus $\bar{x}$]
Upper-case $\bar{X}$ is a random variable --- it describes every possible
value the sample mean could take before data are collected.  Lower-case
$\bar{x}$ is a number, calculated after observing.  Probability statements
are always about $\bar{X}$: writing $P(\bar{x} > 4)$ makes no sense because
$\bar{x}$ is either bigger than~$4$ or it is not.
\end{remark}

\textbf{Mean and variance of the sample mean.}
Because the $X_i$ are independent and identically distributed,
\[
E(\bar{X}) = \mu, \qquad
\mathrm{Var}(\bar{X}) = \frac{\sigma^2}{n}, \qquad
\mathrm{s.d.}(\bar{X}) = \frac{\sigma}{\sqrt{n}}.
\]
The quantity $\sigma/\sqrt{n}$ is called the \textbf{standard error} of the
sample mean.  It measures how much $\bar{X}$ typically varies from sample to
sample: the larger $n$ is, the smaller the standard error.

\subsubsection*{Example: rolling a die three times}
Roll a fair die $3$ times.  Before rolling, $X_1, X_2, X_3$ are each
uniform on $\{1,2,\ldots,6\}$, and $\bar{X}$ is a random variable with
\[
E(\bar{X}) = \frac{7}{2}, \qquad
\mathrm{Var}(\bar{X}) = \frac{35/12}{3} = \frac{35}{36}.
\]
Now roll and get $2, 5, 6$.  Then
$\bar{x} = \frac{2+5+6}{3} = \frac{13}{3} \approx 4.33$, which is one
observed value of the random variable~$\bar{X}$.  Roll three more times
and you will most likely get a different $\bar{x}$.

\subsubsection*{The Central Limit Theorem}

\textbf{Central Limit Theorem.}
Let $X_1, X_2, \ldots, X_n$ be independent observations from \emph{any}
population with mean~$\mu$ and variance~$\sigma^2$.  For large~$n$
(a common guideline is $n \geq 30$),
\[
\bar{X} \approx N\!\left(\mu,\, \frac{\sigma^2}{n}\right).
\]
The population can be discrete, continuous, symmetric, or skewed ---
the distribution of~$\bar{X}$ still approaches a normal distribution as
$n$ increases.  This is why the normal distribution appears so often in
statistics.

\subsubsection*{Special case: Bernoulli population}

If each $X_i$ takes only the values $0$ (failure) and $1$ (success) with
$P(X_i = 1) = p$, then $X_i \sim \mathrm{Bernoulli}(p)$ with
$\mu = p$ and $\sigma^2 = p(1-p)$.  The sample mean becomes
\[
\bar{X} = \frac{1}{n}\sum_{i=1}^{n} X_i = \frac{X}{n} = \hat{p},
\]
the sample proportion.  Substituting into the CLT gives
\[
\hat{p} \approx N\!\left(p,\, \frac{p(1-p)}{n}\right),
\]
which recovers the sample-proportion formulas.  The normal approximation to
the binomial is also a consequence: $X = n\hat{p} \approx N(np,\, np(1-p))$.

\subsubsection*{Why it matters}

In the previous syllabus, the CLT was only applied to Bernoulli populations
(success or failure).  In the current syllabus the CLT is stated for any
population, so questions may involve discrete distributions (such as uniform
dice or scoring rules), continuous probability density functions, or any
other model with a well-defined mean and variance.
